%% intro.tex
%%
%% Copyright 2017 Evandro Coan
%% Copyright 2012-2016 by abnTeX2 group at http://www.abntex.net.br/
%%
%% This work may be distributed and/or modified under the
%% conditions of the LaTeX Project Public License, either version 1.3
%% of this license or (at your option) any later version.
%% The latest version of this license is in
%%   http://www.latex-project.org/lppl.txt
%% and version 1.3 or later is part of all distributions of LaTeX
%% version 2005/12/01 or later.
%%
%% This work has the LPPL maintenance status `maintained'.
%% The Current Maintainer of this work is the Evandro Coan.
%%
%% The last Maintainer of this work was the abnTeX2 team, led
%% by Lauro César Araujo. Further information are available on
%% https://www.abntex.net.br/
%%
%% This work consists of a bunch of files. But originally there ware 3 files
%% which are renamed as follows:
%% Renamed the `abntex2-modelo-include-comandos` to `chapters/chapter_1.tex`
%% Renamed the `abntex2-modelo-trabalho-academico.tex` to `chapters/intro.tex`
%% Renamed the `abntex2-modelo-references.bib` to `aftertext/modelo-ufsc-references.bib`
%%
%% This file was originally the main template file, however this main file was
%% split into several new files, which are respectively drastically changed,
%% except this files which contains most of the main documentation message.
%%

% ------------------------------------------------------------------------
% ------------------------------------------------------------------------
% abnTeX2: Modelo de Trabalho Academico (tese de doutorado, dissertacao de
% mestrado e trabalhos monograficos em geral) em conformidade com
% ABNT NBR 14724:2011: Informacao e documentacao - Trabalhos academicos -
% Apresentacao
% ------------------------------------------------------------------------
% ------------------------------------------------------------------------

% The \phantomsection command is needed to create a link to a place in the document that is not a
% figure, equation, table, section, subsection, chapter, etc.
% https://tex.stackexchange.com/questions/44088/when-do-i-need-to-invoke-phantomsection
\phantomsection

% https://tex.stackexchange.com/questions/5076/is-it-possible-to-keep-my-translation-together-with-original-text
\chapter{\lang{Introduction}{Introdução}}
\phantomsection

\section{Contextualização e Problema}
A evolução dos ecossistemas de Tecnologia da Informação (TI) nas últimas décadas foi marcada por uma mudança de paradigma, saindo de arquiteturas monolíticas e centralizadas para ambientes altamente distribuídos e orientados a serviços \cite{bisht2025generative}. A adoção de microsserviços, balanceadores de carga dinâmicos e arquiteturas baseadas em nuvem (\emph{Cloud-Native}) trouxe inegáveis vantagens em termos de escalabilidade e disponibilidade, mas introduziu uma complexidade operacional sem precedentes. 

Nesses ambientes modernos, incidentes e falhas frequentemente se propagam por múltiplas camadas de infraestrutura, gerando um volume excessivo de alertas e logs de sistemas simultâneos. As abordagens tradicionais de monitoramento, baseadas em regras estáticas e painéis isolados, possuem limitações intrínsecas para lidar com a natureza efêmera e interdependente dessas arquiteturas, resultando no que a literatura costuma chamar de ``fadiga de alertas'' e num elevado Tempo Médio de Resolução (MTTR - \emph{Mean Time to Resolution}) \cite{zhang2024survey, bisht2025generative}. Consequentemente, a Análise de Causa Raiz (RCA - \emph{Root Cause Analysis}), que consiste na identificação do evento subjacente que gerou uma falha sistêmica, torna-se um dos processos mais custosos e demorados para as equipes de Operações de TI e confiabilidade (SRE).

Embora a aplicação de Inteligência Artificial para Operações de TI (AIOps) tenha surgido para automatizar etapas de correlação de eventos e detecção de anomalias, as soluções convencionais de \emph{Machine Learning} muitas vezes atuam como ``caixas pretas'' e requerem extensa engenharia de \emph{features} ou carecem do raciocínio contextual necessário para explicar o incidente de forma acionável para o operador humano.

\section{Justificativa}
Recentemente, a ascensão dos Grandes Modelos de Linguagem (\emph{Large Language Models} - LLMs) abriu novas perspectivas para a evolução das ferramentas de AIOps. Os LLMs possuem notável capacidade de compreensão de linguagem natural, integração de dados multimodais (logs, métricas e \emph{traces}) e raciocínio contextual para inferir relações de dependência entre componentes distribuídos \cite{wang2025integrating, bisht2025generative}. 

Pesquisas recentes têm demonstrado a viabilidade da utilização de agentes baseados em LLMs e técnicas como RAG (\emph{Retrieval-Augmented Generation}) para a síntese de dados de observabilidade, acelerando a resposta a incidentes ao fornecer diagnósticos detalhados e compreensíveis \cite{wang2024rcagent, goel2025accelerating}. No entanto, o desafio persiste no desenvolvimento de integrações eficazes e eficientes entre esses modelos e as populares pilhas de monitoramento \emph{open-source} do mercado, de forma que o contexto sistêmico exato seja injetado nas requisições do modelo para minimizar o risco de alucinações das IAs generativas \cite{li2025large, ahmed2023recommending}.

Diante deste cenário, e dada a necessidade da indústria na redução de tempo e esforço para o RCA de aplicações, justifica-se o desenvolvimento e validação de uma integração entre observabilidade e LLMs em um ambiente real. Além de promover resiliência aos sistemas, a automatização e enriquecimento do RCA em servidores \emph{cloud} pode transformar significativamente os processos de resposta a incidentes de operações nativas em nuvem.

\section{Objetivos}

\subsection{Objetivo Geral}
O objetivo geral deste trabalho é desenvolver e validar uma prova de conceito (PoC) de um sistema de AIOps que integre ferramentas de monitoramento \emph{open-source} a \emph{Large Language Models} (LLMs) para automatizar e enriquecer a Análise de Causa Raiz (RCA) de alertas em Servidores Cloud.

\subsection{Objetivos Específicos}
Para o alcance do objetivo geral proposto, definem-se os seguintes objetivos específicos:
\begin{itemize}
    \item Configurar a pilha de monitoramento baseada em ferramentas \emph{open-source} (Prometheus e família LGTM) em ambiente de servidores \emph{cloud} (AWS).
    \item Investigar e comparar técnicas de integração entre LLMs e dados de observabilidade (como \emph{Few-Shot Prompting}, \emph{Prompt Chaining} ou RAG), além de diferentes provedores e modelos de LLMs, realizando testes de desempenho e acurácia.
    \item Desenvolver um \emph{middleware} (script de automação em Python) capaz de interceptar alertas, extrair métricas de contexto via API e formatar as requisições para o LLM.
    \item Implementar a técnica de injeção de contexto selecionada na etapa de pesquisa no próprio \emph{middleware}, de modo a adaptar as respostas do modelo à infraestrutura local e mitigar alucinações.
    \item Validar a prova de conceito através da simulação de incidentes deliberados, avaliando a utilidade, a precisão e a latência dos diagnósticos gerados pelo sistema desenvolvido.
\end{itemize}


